% =======================
% Section 2: Knowledge & Planning
% =======================
\section{Knowledge \& Planning}

% ---- Slide 8a: RAG — Two-Phase Pipeline ----
\begin{frame}{RAG: Retrieval-Augmented Generation}
\textbf{Goal}: ground model outputs in external knowledge beyond the training cutoff\\[0.4em]
RAG operates in \textbf{two distinct phases} with different parameters and different failure modes:
\vspace{0.4em}
\begin{columns}[T]
\begin{column}{0.48\textwidth}
\begin{block}{Phase 1 — Offline: Knowledge Base Preparation}
\begin{enumerate}
  \item Ingest documents
  \item Chunk text \textit{(strategy, size, overlap)}
  \item Compute embeddings \textit{(model choice)}
  \item Store in vector index \textit{(distance metric)}
\end{enumerate}
\vspace{0.2em}
\scriptsize Decisions made \emph{once}. Heavily influence downstream retrieval quality.
\end{block}
\end{column}
\begin{column}{0.48\textwidth}
\begin{block}{Phase 2 — Online: Retrieval \& Generation}
\begin{enumerate}
  \item Embed user query
  \item Retrieve top-$k$ chunks \textit{(k, threshold, reranking)}
  \item Filter by metadata if needed
  \item Inject context into prompt
  \item Generate grounded answer \textit{(system prompt, model)}
\end{enumerate}
\end{block}
\end{column}
\end{columns}
\vspace{0.4em}
\small \textit{RAG is not one pipeline — it is two pipelines with different owners, timelines, and tuning levers.\\
RAG alone is insufficient when knowledge must be synthesized across many documents or acted upon — that is where agents come in.}
% Speaker notes:
% Offline = data engineering problem; online = runtime problem.
% Different teams often own each phase in production.
\end{frame}

% ---- Slide 8b: RAG Parameters ----
\begin{frame}{RAG Parameters}
\small
\begin{tabular}{llll}
\toprule
\textbf{Phase} & \textbf{Parameter} & \textbf{What it controls} & \textbf{This Project} \\
\midrule
\multirow{4}{*}{\textbf{Offline}}
  & Chunking strategy & How text is split (fixed, semantic, hierarchical) & Fixed-size \\
  & Chunk size        & Amount of context per chunk                       & 512 tokens \\
  & Chunk overlap     & Context preserved across boundaries               & 20\% \\
  & Embedding model   & How text is converted to vectors                  & Titan Embed v2 \\
\midrule
\multirow{5}{*}{\textbf{Online}}
  & Number of results & Chunks returned per search (top-$k$)              & 5 \\
  & Distance metric   & How vectors are compared                          & Cosine \\
  & Reranking         & Cross-encoder rescoring of top-$k$ results        & None \\
  & Metadata filters  & Narrowing search by document or category          & None \\
  & Query formulation & How search queries are constructed                & Model-driven \\
\midrule
\multirow{2}{*}{\textbf{Generation}}
  & System prompt     & Reasoning instructions, output format             & Custom \\
  & Model             & LLM that generates the answer                     & Claude Sonnet 4.5 \\
\bottomrule
\end{tabular}
\vspace{0.4em}

\small \textit{There is no universally correct configuration — parameter choices depend on domain, document type, and task. Reranking is listed even though unused here: its absence is itself a design choice.}
% Speaker notes:
% Walk through each phase group.
% "This Project" column makes concepts concrete.
% Connect each parameter to a real trade-off.
\end{frame}

% ---- Slide 8c: RAG — What Can Go Wrong ----
\begin{frame}{RAG: What Can Go Wrong}
Each parameter is a failure mode waiting to happen:
\vspace{0.3em}
\begin{center}
\small
\begin{tabular}{lll}
\toprule
\textbf{Phase} & \textbf{Bad parameter choice} & \textbf{Failure mode} \\
\midrule
Offline & Chunk size too small          & Loss of context; incoherent retrieval \\
Offline & No chunk overlap              & Answers split across chunk boundaries \\
Offline & Weak embedding model          & Semantically irrelevant chunks retrieved \\
\midrule
Online  & $k$ too low                   & Relevant content not in context \\
Online  & No reranking                  & Top-$k$ by vector similarity $\neq$ top-$k$ by relevance \\
Online  & No metadata filters           & Noise from unrelated documents \\
Online  & Poor query formulation        & Retrieval misses the user's intent \\
\midrule
Generation & Weak system prompt         & Inconsistent format; hallucinated synthesis \\
Generation & Model ignores context      & Generates from parametric memory, not retrieved chunks \\
\bottomrule
\end{tabular}
\end{center}
\vspace{0.3em}
\small \textit{``Stale or irrelevant retrieval'' is one of the most common agent failure modes — it starts here.}
% Speaker notes:
% This table is the payoff: students remember parameters when they understand what breaks.
% Connect to failure modes appendix.
\end{frame}

% ---- Slide 9: Mental Model — Plan Creation then Execution ----
\begin{frame}{My Mental Model: Agents for Plan Creation, Then Execution}
\begin{columns}[T]
\begin{column}{0.52\textwidth}
A practical framing for when to use agents:
\begin{enumerate}
  \item Use an agent to \textbf{create or refine a plan} when the solution path is uncertain
  \item Then \textbf{execute} using either deterministic code or another agent
\end{enumerate}
\vspace{0.4em}
\begin{itemize}
  \item Helps decide whether an agent is actually needed
  \item The boundary is often a gray area
\end{itemize}
\vspace{0.4em}
\textit{``Agents are especially useful when the path to the solution is uncertain, not just when the final answer is hard.''}
\end{column}
\begin{column}{0.44\textwidth}
\begin{tikzpicture}[
  box/.style={draw, rounded corners=4pt, minimum width=3.0cm,
              minimum height=0.6cm, align=center, font=\small},
  every node/.style={font=\small}
]
  \node[box, fill=red!10]    (prob) at (0,3.4) {Uncertain Problem};
  \node[box, fill=orange!15] (plan) at (0,2.4) {Agent creates / refines plan};
  \node[box, fill=blue!10]   (exec) at (0,1.4) {Execution layer};
  \node[box, fill=green!15, minimum width=1.3cm] (det) at (-0.85,0.4) {\scriptsize Deterministic\\\scriptsize Software};
  \node[box, fill=yellow!20, minimum width=1.3cm] (rt)  at ( 0.85,0.4) {\scriptsize Runtime\\\scriptsize Agent};
  \draw[-{Stealth}] (prob) -- (plan);
  \draw[-{Stealth}] (plan) -- (exec);
  \draw[-{Stealth}] (exec) -- (det);
  \draw[-{Stealth}] (exec) -- (rt);
\end{tikzpicture}
\end{column}
\end{columns}
% Speaker notes:
% Examples: "I don't know the schema" → agent discovers.
% "I don't know where info lives" → agent searches.
% Sometimes the best use of an agent is to produce deterministic software.
\end{frame}

% ---- Slide 10: Deterministic Code vs Runtime Agent ----
\begin{frame}{Deterministic Code vs.\ Runtime Agent}
After planning, execution can follow different paths:
\vspace{0.4em}
\begin{columns}[T]
\begin{column}{0.48\textwidth}
\begin{block}{Case 1 — Deterministic Execution}
\begin{itemize}
  \item Plan is well specified
  \item Environment is stable
  \item Limited unforeseen events
  \item Output: scripts, workflows, decision trees
\end{itemize}
\end{block}
\end{column}
\begin{column}{0.48\textwidth}
\begin{block}{Case 2 — Runtime Agentic Execution}
\begin{itemize}
  \item Plan remains high-level
  \item Environment is dynamic / non-stationary
  \item System must adapt step by step
  \item Output: ongoing reasoning + actions
\end{itemize}
\end{block}
\end{column}
\end{columns}
\vspace{0.5em}
\begin{center}
\small
\begin{tabular}{lll}
\toprule
\textbf{Dimension}   & \textbf{Deterministic code} & \textbf{Runtime agent} \\
\midrule
Plan specificity     & High              & Medium / Low \\
Environment          & Stable            & Dynamic \\
Adaptation needed    & Low               & High \\
Cost / latency       & Low, predictable  & Higher, variable \\
Debuggability        & Easy              & Harder \\
\bottomrule
\end{tabular}
\end{center}
\vspace{0.3em}
\textit{``Use software where behavior should be fixed. Use runtime agents where behavior must adapt.''}
% Speaker notes:
% Not every problem should stay agentic at runtime.
% Sometimes the best use of an agent is to produce deterministic software.
\end{frame}
